\section{Formas diferenciales y teoremas integrales}
\label{sec:formas_diferenciales}

\subsection{Teorema de Fubini}
\label{subsec:fubini}

\begin{teorema}[Fubini]
  \label{teo:fubini}
  Sea $f$ una funci\'on acotada en $D=[a,b]\otimes[c,d]$ y supongamos que el
  conjunto de discontinuidades $\Ga$ tiene \'area cero. Si cada l\'inea
  paralela a los ejes coordenados alcanza $\Ga$ en un n\'umero finito de
  puntos, entonces
  \begin{equation}
    \label{eq:fubini}
    \iint_{D}f\dd S
      =\int_{a}^{b}\!\!\int_{c}^{d}f(x,y)\dd y\dd x
      =\int_{c}^{d}\!\!\int_{a}^{b}f(x,y)\dd x\dd y .
  \end{equation}
\end{teorema}

\subsection{\'Algebra exterior}
\label{subsec:algebra_exterior}

\begin{definicion}[Cero-forma]
  \label{def:cero_forma}
  Una \textbf{cero-forma diferencial} es una funci\'on diferenciable
  $f:\R^n\rightarrow\R$. Escribimos $f\in\Ext{0}$, donde $\Ext{0}$ es el
  espacio de funciones sobre el espacio lineal $L$.
\end{definicion}

\begin{definicion}[Mapeo diferencial]
  \label{def:mapeo_diferencial}
  El \textbf{mapeo diferencial} $d$ act\'ua como
  $d:\Ext{n-1}\rightarrow\Ext{n}$. Sobre una cero-forma,
  \begin{equation}
    \label{eq:d_cero_forma}
    \mathbf{d}f=\pdv{f}{x_1}\dd x_1+\pdv{f}{x_2}\dd x_2
      +\dots+\pdv{f}{x_n}\dd x_n\in\Ext{1}.
  \end{equation}
\end{definicion}

\begin{definicion}[$n$-formas]
  \label{def:n_formas}
  Definimos el espacio de $n$-formas diferenciales como
  \begin{equation}
    \label{eq:espacio_n_formas}
    \Ext{n}=\underbrace{\Ext{1}\otimes\Ext{1}\otimes\dots\otimes\Ext{1}}
      _{n\ \text{factores}} ,
  \end{equation}
  cuyos elementos $\om\in\Ext{n}$ son
  \begin{equation}
    \label{eq:n_forma}
    \small
    \om=\sum_{i_1,\dots,i_n}
      \om_{i_1\dots i_n}(x_1,\dots,x_m)
      \dd x_{i_1}\wedge\dd x_{i_2}\wedge\dots\wedge\dd x_{i_n}.
  \end{equation}
\end{definicion}

\subsubsection{El producto cu\~na}
\label{subsubsec:producto_cuna_formas}

\begin{definicion}[Producto cu\~na de formas]
  \label{def:cuna_formas}
  \begin{equation}
    \label{eq:cuna_formas}
    \wedge:\Ext{p}\otimes\Ext{q}\longrightarrow\Ext{p+q},
  \end{equation}
  con las siguientes propiedades, si $\la\in\Ext{p}$ y $\mu\in\Ext{q}$:
  \begin{itemize}
    \item $\la\wedge\mu$ es distributivo;
    \item $\la\wedge(\mu\wedge\nu)=(\la\wedge\mu)\wedge\nu$, es decir, es
          asociativo;
    \item $\mu\wedge\la=(-1)^{pq}\la\wedge\mu$.
  \end{itemize}
\end{definicion}

De ah\'i se siguen las propiedades del \'algebra exterior:
\begin{align}
  \label{eq:exterior_dist_izq}
  (a_1\al_1+a_2\al_2)\wedge\be
    &=a_1\al_1\wedge\be+a_2\al_2\wedge\be,\\
  \label{eq:exterior_dist_der}
  \al\wedge(b_1\be_1+b_2\be_2)
    &=b_1\al\wedge\be_1+b_2\al\wedge\be_2,\\
  \label{eq:exterior_antisim}
  \al\wedge\be&=-\be\wedge\al .
\end{align}
En particular \eqref{eq:exterior_antisim} implica
$\dd x_i\wedge\dd x_i=0$, y por definici\'on $d(\dd x_i)=0$.

\begin{ejemplo}
  \label{ej:cuna_descompuesta}
  Si $\la=\la_1\wedge\dots\wedge\la_p$ y
  $\mu=\mu_1\wedge\dots\wedge\mu_q$, entonces
  \begin{equation}
    \label{eq:cuna_descompuesta}
    \la\wedge\mu
      =\la_1\wedge\dots\wedge\la_p\wedge\mu_1\wedge\dots\wedge\mu_q .
  \end{equation}
\end{ejemplo}

\begin{ejemplo}[Una $1$-forma por una $2$-forma]
  \label{ej:uno_por_dos}
  Sean $\la\in\Ext{1}$ y $\mu\in\Ext{2}$ en $\R^3$,
  \begin{equation}
    \label{eq:lambda_mu}
    \begin{aligned}
      \la&=\la_1\dd x+\la_2\dd y+\la_3\dd z,\\
      \mu&=\mu_1\dd y\wedge\dd z+\mu_2\dd z\wedge\dd x
          +\mu_3\dd x\wedge\dd y .
    \end{aligned}
  \end{equation}
  Todos los productos con un diferencial repetido se anulan y los tres que
  sobreviven se reordenan al mismo signo, de modo que
  \begin{equation}
    \label{eq:lambda_cuna_mu}
    \la\wedge\mu=\qty{\la_1\mu_1+\la_2\mu_2+\la_3\mu_3}
      \dd x\wedge\dd y\wedge\dd z .
  \end{equation}
\end{ejemplo}

\subsubsection{F\'ormula general}
\label{subsubsec:formula_general_cuna}

Para $k$ vectores $\vb{v}_1,\dots,\vb{v}_k\in V$, su producto exterior es
\begin{equation}
  \label{eq:cuna_general}
  \small
  \vb{v}_1\wedge\dots\wedge\vb{v}_k
    =\sum_{1\le i_1<\dots<i_k\le n}
    \begin{vmatrix}
      v_1^{i_1} & \dots & v_1^{i_k}\\
      \vdots & \ddots & \vdots\\
      v_k^{i_1} & \dots & v_k^{i_k}
    \end{vmatrix}
    \vh{e}_{i_1}\wedge\dots\wedge\vh{e}_{i_k}.
\end{equation}

\begin{ejemplo}
  \label{ej:cuna_general_R3}
  Calculemos el producto exterior de dos vectores de $\R^3$ en la base
  can\'onica:
  \begin{equation}
    \label{eq:cuna_general_R3}
    \small
    \begin{aligned}
      \vb{v}\wedge\vb{w}
        &=\sum_{1\le i<j\le 3}
        \begin{vmatrix}
          v_i & v_j\\
          w_i & w_j
        \end{vmatrix}
        \vh{e}_i\wedge\vh{e}_j\\
        &=(v_1w_2-v_2w_1)\,\vh{e}_1\wedge\vh{e}_2\\
        &\quad+(v_1w_3-v_3w_1)\,\vh{e}_1\wedge\vh{e}_3\\
        &\quad+(v_2w_3-v_3w_2)\,\vh{e}_2\wedge\vh{e}_3 ,
    \end{aligned}
  \end{equation}
  donde $\vh{e}_i\wedge\vh{e}_j\equiv\vh{e}_{ij}$. Identificando
  $\vh{e}_{23}\mapsto\vh{e}_1$, $\vh{e}_{31}\mapsto\vh{e}_2$ y
  $\vh{e}_{12}\mapsto\vh{e}_3$ se recupera el producto cruz
  \eqref{eq:cuna_3d}.
\end{ejemplo}

\subsection{La derivada exterior en $\R^3$}
\label{subsec:derivada_exterior}

Sea $\om\in\Ext{1}$ con $\om_i:\R^3\rightarrow\R$,
\begin{equation}
  \label{eq:uno_forma}
  \om=\om_1\dd x+\om_2\dd y+\om_3\dd z .
\end{equation}
Entonces, usando $d(\dd x_i)=0$,
\begin{equation}
  \label{eq:d_omega_paso1}
  d\om=d\om_1\wedge\dd x+d\om_2\wedge\dd y+d\om_3\wedge\dd z,
\end{equation}
y desarrollando cada $d\om_i$ con \eqref{eq:d_cero_forma},
\begin{equation}
  \label{eq:d_omega_paso2}
  \small
  \begin{aligned}
    d\om&=\qty{\pdv{\om_1}{x}\dd x+\pdv{\om_1}{y}\dd y
              +\pdv{\om_1}{z}\dd z}\wedge\dd x\\
        &\quad+\qty{\pdv{\om_2}{x}\dd x+\pdv{\om_2}{y}\dd y
              +\pdv{\om_2}{z}\dd z}\wedge\dd y\\
        &\quad+\qty{\pdv{\om_3}{x}\dd x+\pdv{\om_3}{y}\dd y
              +\pdv{\om_3}{z}\dd z}\wedge\dd z .
  \end{aligned}
\end{equation}
Los t\'erminos con diferencial repetido se anulan y los restantes se agrupan:
\begin{equation}
  \label{eq:d_omega_rot}
  \begin{aligned}
    d\om&=\qty{\pdv{\om_2}{x}-\pdv{\om_1}{y}}\dd x\wedge\dd y\\
        &\quad+\qty{\pdv{\om_1}{z}-\pdv{\om_3}{x}}\dd z\wedge\dd x\\
        &\quad+\qty{\pdv{\om_3}{y}-\pdv{\om_2}{z}}\dd y\wedge\dd z .
  \end{aligned}
\end{equation}

Ahora bien, el diferencial de superficie se descompone como
\begin{equation}
  \label{eq:dS_componentes}
  \left\{
  \begin{aligned}
    \dd\vb{S}_x&=\vh{x}\dd y\wedge\dd z,\\
    \dd\vb{S}_y&=\vh{y}\dd z\wedge\dd x,\\
    \dd\vb{S}_z&=\vh{z}\dd x\wedge\dd y,
  \end{aligned}
  \right.
\end{equation}
de modo que \eqref{eq:d_omega_rot} se lee como un producto punto:
\begin{equation}
  \label{eq:d_omega_producto}
  d\om=
  \underbrace{\qty{\nabla\wedge
    \begin{pmatrix}\om_1\\\om_2\\\om_3\end{pmatrix}}}_{\rot\boldsymbol{\om}}
  \cdot
  \underbrace{
  \begin{pmatrix}
    \dd y\wedge\dd z\\
    \dd z\wedge\dd x\\
    \dd x\wedge\dd y
  \end{pmatrix}}_{\dd\vb{S}} .
\end{equation}

\subsection{Teorema de Stokes}
\label{subsec:stokes}

\begin{teorema}[Stokes generalizado]
  \label{teo:stokes_general}
  \begin{equation}
    \label{eq:stokes_general}
    \int_{\partial S}\om=\int_{S}d\om .
  \end{equation}
\end{teorema}

Recordando que para $\om\in\Ext{1}$ se tiene
$\om=\boldsymbol{\om}\cdot\dd\vb{r}$, la combinaci\'on de
\eqref{eq:stokes_general} con \eqref{eq:d_omega_producto} da la forma
vectorial cl\'asica
\begin{equation}
  \label{eq:stokes_vectorial}
  \int_{C=\partial S}\vb{F}\cdot\dd\vb{r}
    =\iint_{S}\rot\vb{F}\cdot\dd\vb{S} .
\end{equation}

\begin{observacion}
  El teorema s\'olo es v\'alido en superficies \emph{abiertas}: una superficie
  cerrada no tiene frontera, de modo que el miembro izquierdo no existe.
\end{observacion}

\begin{observacion}[Invariancia de norma]
  Supongamos $\vb{G}=\rot\vb{F}$. Como el rotacional de un gradiente es cero,
  tambi\'en $\vb{G}=\rot\qty{\vb{F}+\nabla g}$ para cualquier campo escalar
  $g$ de clase $C^2$, y entonces
  \begin{equation}
    \label{eq:invariancia_norma}
    \iint_{S}\vb{G}\cdot\dd\vb{S}
      =\int_{\partial S}\qty{\vb{F}+\nabla g}\cdot\dd\vb{r},
  \end{equation}
  es decir, $\vb{F}$ est\'a determinado s\'olo salvo un gradiente.
\end{observacion}

\subsubsection{Ley de Amp\`ere}
\label{subsubsec:ampere}

En su forma integral,
\begin{equation}
  \label{eq:ampere_integral}
  \int_{C=\partial S}\vb{B}\cdot\dd\vb{r}=\mu_0 I .
\end{equation}
Escribiendo la corriente como el flujo de la densidad de corriente
$\vb{J}$ y aplicando \eqref{eq:stokes_vectorial} al miembro izquierdo,
\begin{equation}
  \label{eq:ampere_paso}
  \iint_{S}\rot\vb{B}\cdot\dd\vb{S}
    =\mu_0\iint_{S}\vb{J}\cdot\dd\vb{S},
\end{equation}
de donde
\begin{equation}
  \label{eq:ampere_diferencial}
  \iint_{S}\qty{\rot\vb{B}-\mu_0\vb{J}}\cdot\dd\vb{S}=0
  \quad\forall S .
\end{equation}
Como esto vale para \emph{toda} superficie, el integrando debe anularse:
\begin{equation}
  \label{eq:ampere_local}
  \rot\vb{B}=\mu_0\vb{J},
  \qquad\text{es decir}\qquad
  \nabla\wedge\vb{B}=\mu_0\vb{J} .
\end{equation}
Aqu\'i $\vb{B}=\rot\vb{A}$, con $\vb{A}$ el \textbf{potencial vectorial};
por la observaci\'on anterior, $\vb{A}$ y $\vb{A}+\nabla\la$ producen el mismo
campo, que es la libertad de norma del electromagnetismo.

\subsection{Teorema de Gauss}
\label{subsec:gauss}

Sea ahora $\om\in\Ext{2}$,
\begin{equation}
  \label{eq:dos_forma}
  \om=\om_1\dd y\wedge\dd z+\om_2\dd z\wedge\dd x
    +\om_3\dd x\wedge\dd y .
\end{equation}
Repitiendo el c\'alculo de \eqref{eq:d_omega_paso2}, ahora sobreviven
\'unicamente los t\'erminos que completan las tres coordenadas:
\begin{equation}
  \label{eq:d_omega_div}
  \begin{aligned}
    d\om&=\qty{\pdv{\om_1}{x}+\pdv{\om_2}{y}+\pdv{\om_3}{z}}
          \dd x\wedge\dd y\wedge\dd z\\
        &=\qty{\nabla\cdot\boldsymbol{\om}}\dd V .
  \end{aligned}
\end{equation}

\begin{teorema}[Divergencia]
  \label{teo:divergencia}
  \begin{equation}
    \label{eq:divergencia}
    \iint_{S=\partial V}\vb{F}\cdot\dd\vb{S}
      =\iiint_{V}\dive\vb{F}\dd V .
  \end{equation}
  Si la superficie no es cerrada, el teorema no aplica.
\end{teorema}

\begin{ejemplo}
  \label{ej:divergencia}
  Sea
  \begin{equation}
    \label{eq:F_divergencia}
    \vb{F}=
    \begin{pmatrix}
      3xy^2\\3yx^2\\z^3
    \end{pmatrix}
  \end{equation}
  y busquemos su flujo sobre la esfera unitaria. Al ser una superficie cerrada
  conviene el teorema de la divergencia:
  \begin{equation}
    \label{eq:div_F}
    \dive\vb{F}=\nabla\cdot\vb{F}=3y^2+3x^2+3z^2=3\qty{x^2+y^2+z^2},
  \end{equation}
  que en esf\'ericas es simplemente $\dive\vb{F}=3r^2$. Entonces
  \begin{equation}
    \label{eq:flujo_esfera}
    \begin{aligned}
      \iint_{S=\partial V}\vb{F}\cdot\dd\vb{S}
        &=\int_{0}^{2\pi}\!\!\int_{0}^{\pi}\!\!\int_{0}^{1}
          \qty{3r^2}r^2\sin\theta\dd r\dd\theta\dd\vph\\
        &=3(4\pi)\int_{0}^{1}r^4\dd r=\frac{12}{5}\pi .
    \end{aligned}
  \end{equation}
\end{ejemplo}
